% ============================================================
% paper48_additive_number_theory_en.tex
% Paper 48: Additive Number Theory
% Author: Michael Fuhrmann
% Project: specialist-maths, Build 137
% Date: 2026-03-12
% ============================================================

\documentclass[12pt,a4paper]{amsart}
\usepackage{amsmath,amssymb,amsthm}
\usepackage{mathtools}
\usepackage[utf8]{inputenc}
\usepackage[T1]{fontenc}
\usepackage{hyperref}
\usepackage{enumitem,booktabs,array}
\usepackage{geometry}
\geometry{margin=2.5cm}

% --- Theorem-Umgebungen ---
\newtheorem{theorem}{Theorem}[section]
\newtheorem{lemma}[theorem]{Lemma}
\newtheorem{corollary}[theorem]{Corollary}
\newtheorem{proposition}[theorem]{Proposition}
\newtheorem{conjecture}[theorem]{Conjecture}
\theoremstyle{definition}
\newtheorem{definition}[theorem]{Definition}
\newtheorem{example}[theorem]{Example}
\newtheorem{remark}[theorem]{Remark}

% --- Metadaten ---
\title{Additive Number Theory:\\
Sumsets, Goldbach, Waring, and Beyond}

\author{Michael Fuhrmann}
\address{Specialist Mathematics Project, Build 137}
\email{specialist-maths@localhost}
\date{\today}

% --- Makros ---
\newcommand{\N}{\mathbb{N}}
\newcommand{\Z}{\mathbb{Z}}
\newcommand{\Q}{\mathbb{Q}}
\newcommand{\R}{\mathbb{R}}
\newcommand{\C}{\mathbb{C}}
\newcommand{\Primes}{\mathbb{P}}
\newcommand{\abs}[1]{\left|#1\right|}
\newcommand{\floor}[1]{\left\lfloor #1 \right\rfloor}
\newcommand{\sumset}[2]{#1 + #2}
\newcommand{\ksum}[2]{#1 \hat{+} #2}
\DeclareMathOperator{\ord}{ord}
\DeclareMathOperator{\lcm}{lcm}

% ============================================================
\begin{document}

\begin{abstract}
Additive number theory studies the additive structure of the integers,
asking which numbers can be represented as sums of elements from
prescribed sets---primes, perfect powers, or arbitrary subsets of~$\N$.
This paper provides a self-contained introduction to the central concepts
and results: sumset theory, Schnirelmann density, bases of finite order,
the Goldbach conjectures (binary and ternary), Chen's theorem, the
Hardy--Littlewood heuristic, Waring's problem and the Hilbert--Waring
theorem, Schnirelmann's theorem, Freiman's structure theorem, and the
Erd\H{o}s--Tur\'{a}n conjecture.  We distinguish carefully between
proved theorems and open conjectures, and illustrate all concepts with
computational examples drawn from the accompanying Python modules
\texttt{additive\_number\_theory.py} and \texttt{goldbach\_extended.py}.
\end{abstract}

\maketitle
\tableofcontents

% ============================================================
\section{Introduction}\label{sec:intro}
% ============================================================

Additive number theory is the branch of mathematics that investigates
how integers can be expressed as sums of elements from specific sets.
The simplest question is: given a set $A \subseteq \N$, which positive
integers belong to the \emph{sumset} $A + A = \{a + a' : a, a' \in A\}$?
More generally, one asks for the \emph{$h$-fold sumset}
$hA = A + A + \cdots + A$ ($h$ summands).

Historically, additive number theory grew from three celebrated problems:

\begin{enumerate}[label=(\arabic*)]
\item \textbf{Goldbach (1742):} Is every even integer $\geq 4$ the sum
      of two primes?  Is every odd integer $\geq 7$ the sum of three primes?
\item \textbf{Waring (1770):} Can every positive integer be expressed as
      a sum of at most $g(k)$ perfect $k$-th powers?
\item \textbf{Basis problems:} Which sets $A$ satisfy $hA = \N$ for
      some finite $h$?  What is the minimal such $h$?
\end{enumerate}

The subject has deep connections to analytic number theory
(the Hardy--Littlewood circle method), harmonic analysis (Fourier
techniques on $\Z/N\Z$), combinatorics (Freiman's theorem, additive
combinatorics), and modern arithmetic geometry.

\subsection*{Overview of this paper}

Section~\ref{sec:sumsets} introduces sumsets, Schnirelmann density, and
bases of finite order.  Section~\ref{sec:goldbach} presents Goldbach's
binary and ternary conjectures with their current status.
Section~\ref{sec:goldbach-ext} covers Chen's theorem, Vinogradov's
three-prime theorem, and the Hardy--Littlewood conjecture.
Section~\ref{sec:waring} treats Waring's problem and the Hilbert--Waring
theorem.  Section~\ref{sec:schnirelmann} develops Schnirelmann's density
approach.  Section~\ref{sec:freiman} presents Freiman's structure theorem.
Section~\ref{sec:erdos-turan} states the Erd\H{o}s--Tur\'{a}n conjecture.
Section~\ref{sec:computational} illustrates computational aspects.
Section~\ref{sec:references} collects references.

% ============================================================
\section{Sumsets and Density}\label{sec:sumsets}
% ============================================================

\subsection{Basic Definitions}

\begin{definition}[Sumset]
Let $A, B \subseteq \Z$.  The \emph{sumset} of $A$ and $B$ is
\[
  A + B \;:=\; \{ a + b : a \in A,\; b \in B \}.
\]
The \emph{$h$-fold sumset} is $hA := A + A + \cdots + A$ ($h$ terms).
The \emph{restricted sumset} is
$A \hat{+} B := \{a + b : a \in A,\, b \in B,\, a \neq b\}$.
\end{definition}

\begin{definition}[Basis of Order $h$]\label{def:basis}
A set $A \subseteq \N_0$ is called an \emph{additive basis of order~$h$}
(or \emph{$h$-basis}) for $\N$ if every sufficiently large positive
integer belongs to $hA$.  It is a \emph{basis of exact order $h$} if
$h$ is the smallest such integer.  If only finitely many positive
integers are missing from $hA$, we speak of an \emph{asymptotic basis}.
\end{definition}

The prime numbers, the squares, the cubes, and many other classical sets
form bases of finite order (cf.\ Sections~\ref{sec:goldbach}
and~\ref{sec:waring}).

\subsection{Schnirelmann Density}

The concept of \emph{Schnirelmann density} provides a quantitative
measure for how densely a set is distributed in~$\N$.

\begin{definition}[Schnirelmann Density]\label{def:schnirelmann-density}
For a set $A \subseteq \N_0$ with $0 \in A$, the \emph{Schnirelmann density}
is
\[
  \sigma(A) \;:=\; \inf_{n \geq 1} \frac{A(n)}{n},
\]
where $A(n) = \abs{\{a \in A : 1 \leq a \leq n\}}$ is the \emph{counting
function} of~$A$.
\end{definition}

Note that $\sigma(A) = 0$ if $1 \notin A$, and $\sigma(A) = 1$ if and
only if $A = \N$.  The key advantage over natural density is:

\begin{theorem}[Schnirelmann's Sum Theorem]\label{thm:schnirelmann-sum}
For sets $A, B \subseteq \N_0$ with $0 \in A \cap B$,
\[
  \sigma(A + B) \;\geq\; \sigma(A) + \sigma(B) - \sigma(A)\,\sigma(B).
\]
In particular, if $\sigma(A) + \sigma(B) \geq 1$, then $A + B = \N$.
\end{theorem}

\begin{proof}
Let $n \geq 1$ be arbitrary.  Among the $n + 1$ elements
$0 = a_0 < a_1 < \cdots < a_k \leq n$ of $A \cap [0,n]$, the
intervals $(a_{i-1}, a_i)$ contain no element of $A$, hence the elements
$a_{i-1} + b$ for $b \in B \cap [0, a_i - a_{i-1}]$ all lie in $A + B$.
Counting carefully yields the stated inequality.
\end{proof}

\begin{corollary}\label{cor:density-one}
If $\sigma(A) \geq 1/2$, then $A + A \supseteq \N$.
\end{corollary}

\subsection{Natural Density and Upper Density}

For comparison, recall that the \emph{natural (asymptotic) density} of
$A \subseteq \N$ is $d(A) = \lim_{n \to \infty} A(n)/n$ (if the limit
exists), and the \emph{upper density} is
$\bar d(A) = \limsup_{n \to \infty} A(n)/n$.

Unlike Schnirelmann density, natural density is not sub-multiplicative
and does not directly control the order of~$A$ as a basis.  The
Schnirelmann density is smaller: $\sigma(A) \leq d(A)$ whenever both
exist.

% ============================================================
\section{Goldbach's Conjecture}\label{sec:goldbach}
% ============================================================

\subsection{Historical Background}

In a letter to Leonhard Euler dated 7 June 1742, Christian Goldbach
stated that every integer greater than 2 is the sum of three primes
(counting 1 as prime, following the convention of his time).  Euler
reformulated the conjecture in the modern form, leading to two distinct
statements that we call the \emph{binary} (strong) and \emph{ternary}
(weak) Goldbach conjectures.

\subsection{The Binary (Strong) Goldbach Conjecture}

\begin{conjecture}[Binary Goldbach Conjecture]\label{conj:goldbach-binary}
Every even integer $n \geq 4$ is the sum of two prime numbers:
\[
  \forall\; n \geq 4,\; n \equiv 0 \pmod{2} \;\Longrightarrow\;
  \exists\; p, q \in \Primes : n = p + q.
\]
\end{conjecture}

This conjecture has been open since 1742.  Define the
\emph{Goldbach partition function}
\[
  G(n) \;:=\; \abs{\bigl\{(p, q) : p + q = n,\; p \leq q,\;
                p, q \text{ prime}\bigr\}}.
\]

\begin{remark}
Conjecture~\ref{conj:goldbach-binary} is equivalent to $G(n) \geq 1$
for all even $n \geq 4$.  Numerical evidence is overwhelming:
\begin{itemize}
\item Verified for all even $n \leq 4 \times 10^{18}$
      (Oliveira e Silva, Herzog, Pardi, 2012 \cite{OliveiraSilva2012}).
\item For ``large'' $n$, $G(n)$ grows approximately as
      $n / (\log n)^2$ (see Section~\ref{sec:goldbach-ext}).
\item $G(4) = 1$: $4 = 2 + 2$.
\item $G(100) = 6$: $100 = 3+97 = 11+89 = 17+83 = 29+71 = 41+59 = 47+53$.
\end{itemize}
\end{remark}

\subsection{The Ternary (Weak) Goldbach Conjecture}

\begin{conjecture}[Ternary Goldbach Conjecture --- Now a Theorem]\label{conj:goldbach-ternary}
Every odd integer $n \geq 7$ is the sum of three prime numbers.
\end{conjecture}

We note that this is \emph{no longer a conjecture}: it was proved by
Harald Helf\-gott in 2013, completing the work started by Vinogradov in
1937.

\begin{theorem}[Helfgott 2013, \cite{Helfgott2013}]\label{thm:helfgott}
Every odd integer $n \geq 7$ is expressible as the sum of three primes:
\[
  \forall\; n \geq 7,\; n \equiv 1 \pmod{2} \;\Longrightarrow\;
  \exists\; p_1, p_2, p_3 \in \Primes : n = p_1 + p_2 + p_3.
\]
\end{theorem}

\begin{proof}[Proof sketch]
The proof combines:
\begin{enumerate}[label=(\roman*)]
\item \textbf{Analytic part:} The circle method of Hardy--Littlewood--
      Vinogradov yields an asymptotic formula for the number of
      representations $r_3(n)$ when $n$ is large.  Specifically,
      Vinogradov (1937) showed $r_3(n) > 0$ for all odd $n > \exp(\exp(11.503))$.
      Helfgott improved this to all odd $n > 10^{30}$ via refined
      exponential sum estimates.
\item \textbf{Computational part:} Direct verification for all odd
      $7 \leq n \leq 8.875 \times 10^{30}$ using the classification of
      primes and modern computer algebra.
\end{enumerate}
Together these two parts cover all odd $n \geq 7$.
\end{proof}

\subsection{Heuristic Evidence for the Binary Conjecture}

The \emph{Goldbach comet} is the scatter plot of $G(n)$ for even $n$.
It shows a clear upward trend, with characteristic spikes at
highly composite numbers (cf.\ Section~\ref{sec:computational}).

The heuristic argument is simple: if primes were distributed randomly
with density $1/\log n$ near $n$, then
\[
  G(n) \;\approx\; \frac{n}{(\log n)^2}
\]
which grows without bound.  More precisely, a probabilistic model gives
\[
  G(n) \;\approx\; 2 \Pi_2
  \prod_{\substack{p > 2 \\ p \mid n}} \frac{p-1}{p-2}
  \cdot \frac{n}{(\log n)^2}
\]
where $\Pi_2 = \prod_{p > 2} \frac{p(p-2)}{(p-1)^2} \approx 0.6602$
is the twin prime constant (Hardy--Littlewood conjecture, see
Section~\ref{sec:goldbach-ext}).

% ============================================================
\section{Goldbach Extended: Chen, Vinogradov, Hardy--Littlewood}
\label{sec:goldbach-ext}
% ============================================================

\subsection{Chen's Theorem}

The strongest unconditional approximation to the binary Goldbach
conjecture is:

\begin{theorem}[Chen's Theorem, 1973 \cite{Chen1973}]\label{thm:chen}
Every sufficiently large even integer $n$ can be written as
\[
  n = p + m,
\]
where $p$ is prime and $m$ is either prime or a product of exactly two
primes (i.e., $m \in P_2$, a \emph{$P_2$-number}).
\end{theorem}

\begin{remark}
Chen's theorem is often stated as ``$p + P_2$''.  It falls just short of
Goldbach: instead of requiring $m$ to be prime, it allows $m = p_1 p_2$
as well.  The proof uses the \emph{large sieve} and sophisticated sieve
methods.
\end{remark}

\begin{example}
For $n = 20$:
\begin{align*}
  20 &= 3 + 17 \quad (17 \text{ prime}) \\
  20 &= 7 + 13 \quad (13 \text{ prime}) \\
  20 &= 2 + 18 = 2 + 2 \cdot 9 \quad (18 = 2 \cdot 3^2 \text{: not } P_2) \\
  20 &= 7 + 13 \checkmark
\end{align*}
Chen's theorem guarantees at least one such representation for all large
even $n$.
\end{example}

\subsection{Vinogradov's Three-Prime Theorem}

\begin{theorem}[Vinogradov 1937 \cite{Vinogradov1937}]\label{thm:vinogradov}
Every sufficiently large odd integer $N$ is the sum of three primes:
\[
  r_3(N) \;:=\; \sum_{\substack{p_1 + p_2 + p_3 = N \\ p_i \text{ prime}}}
  \log p_1 \cdot \log p_2 \cdot \log p_3
  \;\sim\;
  \mathfrak{S}(N) \cdot \frac{N^2}{2(\log N)^3},
\]
where
\[
  \mathfrak{S}(N) \;=\;
  \prod_{p \mid N} \left(1 - \frac{1}{(p-1)^2}\right)
  \prod_{p \nmid N} \left(1 + \frac{1}{(p-1)^3}\right)
\]
is the \emph{singular series}, which satisfies $\mathfrak{S}(N) > 0$
for all odd~$N$.
\end{theorem}

\begin{proof}[Proof sketch --- Circle Method]
The proof uses the Hardy--Littlewood--Ramanujan circle method.  Define
the exponential sum
\[
  S(\alpha) \;:=\; \sum_{p \leq N} \Lambda(p)\, e^{2\pi i \alpha p},
\]
where $\Lambda$ is the von Mangoldt function.  Then by Fourier analysis,
\[
  r_3(N) \;=\; \int_0^1 S(\alpha)^3\, e^{-2\pi i N \alpha}\, d\alpha.
\]

\textbf{Major arcs} $\mathfrak{M}$: For $\alpha$ near a rational $a/q$
with small denominator $q \leq Q = N^{1/3}$, i.e.,
$\abs{\alpha - a/q} < N^{-2/3}$, one approximates $S(\alpha)$ via
Gauss sums and the prime number theorem in arithmetic progressions.
This yields the main term $\mathfrak{S}(N) \cdot N^2 / (2(\log N)^3)$.

\textbf{Minor arcs} $\mathfrak{m} = [0,1] \setminus \mathfrak{M}$:
Vinogradov's mean value theorem and the method of exponential sums
(van der Corput--Weyl--Vinogradov) give the bound
\[
  \sup_{\alpha \in \mathfrak{m}} |S(\alpha)| \;\ll\; \frac{N}{(\log N)^A}
\]
for every $A > 0$.  Hence the minor arc integral is
$o\!\left(N^2/(\log N)^3\right)$ and negligible.
\end{proof}

\subsection{Hardy--Littlewood Goldbach Conjecture}

\begin{conjecture}[Hardy--Littlewood, 1923 \cite{HardyLittlewood1923}]
\label{conj:hardy-littlewood}
For every even $n \geq 4$,
\[
  G(n) \;\sim\; 2\,\Pi_2
  \prod_{\substack{p > 2 \\ p \mid n}}
  \frac{p-1}{p-2}
  \cdot \frac{n}{(\log n)^2}, \qquad n \to \infty,
\]
where
\[
  \Pi_2 \;:=\; \prod_{p > 2} \frac{p(p-2)}{(p-1)^2}
        \;\approx\; 0.6601618\ldots
\]
is the \emph{twin prime constant}.
\end{conjecture}

\begin{remark}
This conjecture is open.  It is consistent with all numerical evidence.
Note that the extra factor $\prod_{p \mid n} (p-1)/(p-2)$ explains why
$G(n)$ is larger for numbers $n$ with many odd prime factors (the comet
spikes).
\end{remark}

% ============================================================
\section{Waring's Problem}\label{sec:waring}
% ============================================================

\subsection{Statement and History}

In 1770, Edward Waring asserted without proof that every positive integer
is the sum of at most 4 squares, 9 cubes, 19 fourth powers, and so on.
This became known as \emph{Waring's problem}.

\begin{definition}[$g(k)$ and $G(k)$]
For $k \geq 2$, let
\begin{align*}
  g(k) &:= \text{smallest } s \text{ such that every positive integer is
           a sum of } s \text{ perfect } k\text{-th powers,}\\
  G(k) &:= \text{smallest } s \text{ such that every \emph{sufficiently large}
           positive integer is a sum of } s \text{ perfect } k\text{-th powers.}
\end{align*}
Clearly $G(k) \leq g(k)$.
\end{definition}

\subsection{The Hilbert--Waring Theorem}

\begin{theorem}[Hilbert--Waring Theorem, 1909 \cite{Hilbert1909}]
\label{thm:hilbert-waring}
For every $k \geq 2$, the quantity $g(k)$ is finite.  Equivalently,
the set of perfect $k$-th powers forms an additive basis of finite order.
\end{theorem}

\begin{proof}[Proof idea]
Hilbert's original proof (1909) used algebraic identities and a
combinatorial averaging argument.  Modern proofs use the Hardy--Littlewood
circle method: the number of representations of $n$ as a sum of $s$
perfect $k$-th powers is
\[
  R_{s,k}(n) \;=\; \int_0^1
  \left(\sum_{m=1}^{n^{1/k}} e^{2\pi i \alpha m^k}\right)^s
  e^{-2\pi i n \alpha}\, d\alpha.
\]
For $s$ sufficiently large (relative to $k$), the method shows
$R_{s,k}(n) > 0$ for all large $n$.
\end{proof}

\subsection{Known Values}

The exact value of $g(k)$ is known for all $k$ (with minor exceptions
depending on a Diophantine condition):

\begin{table}[h]
\centering
\begin{tabular}{@{}ccc@{}}
\toprule
$k$ & $g(k)$ & $G(k)$ (known/conjectured) \\
\midrule
2 & 4  & 4 \\
3 & 9  & $\leq 7$ (exact: $G(3) = 7$ conjectured) \\
4 & 19 & 15 (conjectured; $G(4) \leq 15$ proved) \\
5 & 37 & $\leq 21$ \\
6 & 73 & $\leq 31$ \\
7 & 143 & $\leq 45$ \\
8 & 279 & $\leq 62$ \\
\bottomrule
\end{tabular}
\caption{Values of $g(k)$ (all proved) and bounds on $G(k)$.}
\label{tab:waring}
\end{table}

\begin{theorem}[Lagrange's Four-Square Theorem, 1770 \cite{Lagrange1770}]
Every positive integer is the sum of four perfect squares:
\[
  g(2) = 4.
\]
\end{theorem}

The general formula for $g(k)$ (for $k \geq 2$) is:
\[
  g(k) = 2^k + \floor{(3/2)^k} - 2
\]
whenever $2^k\{(3/2)^k\} + \floor{(3/2)^k} \leq 2^k$
(this holds for all known $k$, and is expected to hold for all $k$).

\begin{example}
\begin{align*}
  g(2) &= 4: \quad 7 = 4 + 1 + 1 + 1 = 2^2 + 1^2 + 1^2 + 1^2. \\
  g(3) &= 9: \quad 23 = 8 + 8 + 1 + 1 + 1 + 1 + 1 + 1 + 1 = 9 \text{ cubes}. \\
  g(4) &= 19: \quad 15 = 1^4 + \cdots + 1^4 \text{ (15 times)}.
\end{align*}
\end{example}

% ============================================================
\section{Schnirelmann's Theorem}\label{sec:schnirelmann}
% ============================================================

\subsection{The Schnirelmann Approach to Goldbach}

Schnirelmann (1930) introduced his density to give the first unconditional
proof that the primes form an additive basis of \emph{finite} order.

\begin{theorem}[Schnirelmann, 1930 \cite{Schnirelmann1930}]
\label{thm:schnirelmann}
The set of primes $\mathbb{P}$ (together with $\{0, 1\}$) forms an
additive basis of order $s < \infty$ for $\N$, i.e., every integer
$n \geq 2$ is the sum of at most $s$ primes.
\end{theorem}

\begin{proof}[Proof sketch]
Let $A = \{0, 1\} \cup \mathbb{P}$.  Goldbach's conjecture (at the time
unproved) would give $\sigma(A + A) > 0$.  Schnirelmann cleverly avoided
Goldbach by using:
\begin{enumerate}[label=(\arabic*)]
\item $\sigma(A) > 0$ (since $1 \in A$ and many primes are in $A$).
\item The \emph{Schnirelmann sum theorem}
      (Theorem~\ref{thm:schnirelmann-sum}): iterating $A + A + \cdots$
      eventually has Schnirelmann density $\geq 1$, i.e., equals all of $\N$.
\end{enumerate}
Schnirelmann's original bound was $s \leq 800{,}000$.  Later improvements
(Romanov, Heilbronn, Landau, Vinogradov) reduced this dramatically.
\end{proof}

\subsection{The Schnirelmann Constant}

\begin{definition}[Schnirelmann Constant]
The \emph{Schnirelmann constant} $s$ is the smallest integer such that
every integer $n \geq 2$ is the sum of at most $s$ primes.
\end{definition}

\begin{table}[h]
\centering
\begin{tabular}{@{}lll@{}}
\toprule
Author & Year & Bound on $s$ \\
\midrule
Schnirelmann & 1930 & $s \leq 800{,}000$ \\
Romanov      & 1934 & $s \leq 2{,}208$ \\
Vinogradov   & 1937 & $s \leq 4$ (for odd $n \geq N_0$ large) \\
Helfgott     & 2013 & $s \leq 4$ (for all $n \geq 2$; $s = 3$ for odd $n \geq 7$) \\
\bottomrule
\end{tabular}
\caption{History of bounds on the Schnirelmann constant~$s$.}
\label{tab:schnirelmann}
\end{table}

As a consequence of Helfgott's theorem (Theorem~\ref{thm:helfgott}):
\begin{corollary}
Every integer $n \geq 2$ is the sum of at most $4$ primes.
\end{corollary}

\subsection{Schnirelmann Density Inequality}

For the set of primes $\mathbb{P}$, the Schnirelmann density satisfies:
\[
  \sigma(\mathbb{P}) \;=\; 0
\]
since $2 \notin \mathbb{P}_0$ (where we adjoin $0$).  Nevertheless,
the density of $A = \{0,1\} \cup \mathbb{P}$ satisfies $\sigma(A) > 0$
because $1 \in A$.  The key inequality
\[
  \sigma(A + B) \;\geq\; \sigma(A) + \sigma(B) - \sigma(A)\sigma(B)
\]
allows iterative improvement until density reaches~1.

% ============================================================
\section{Freiman's Theorem}\label{sec:freiman}
% ============================================================

\subsection{The Trivial Bound}

\begin{lemma}[Trivial Lower Bound]\label{lem:trivial}
For any finite set $A \subseteq \Z$ with $\abs{A} = n$,
\[
  \abs{A + A} \;\geq\; 2n - 1.
\]
\end{lemma}

\begin{proof}
If $A = \{a_1 < a_2 < \cdots < a_n\}$, then
$2a_1 < a_1 + a_2 < 2a_2 < \cdots < a_{n-1} + a_n < 2a_n$
gives $2n - 1$ distinct elements.
\end{proof}

Equality $\abs{A + A} = 2\abs{A} - 1$ holds if and only if $A$ is an
arithmetic progression (AP).  Freiman's theorem gives a structural
characterization when $\abs{A + A}$ is small.

\subsection{Freiman's Structure Theorem}

\begin{definition}[Generalized Arithmetic Progression]
A \emph{generalized arithmetic progression} (GAP) of dimension $d$ is
a set of the form
\[
  P \;=\; \{a_0 + x_1 r_1 + x_2 r_2 + \cdots + x_d r_d :
             0 \leq x_j \leq L_j - 1\}
\]
for some $a_0, r_1, \ldots, r_d \in \Z$ and positive integers
$L_1, \ldots, L_d$.  Its \emph{size} is $L_1 \cdot L_2 \cdots L_d$.
\end{definition}

\begin{theorem}[Freiman's Theorem \cite{Freiman1973}]\label{thm:freiman}
Let $A \subseteq \Z$ be a finite set with $\abs{A + A} \leq K \abs{A}$
for some $K \geq 1$.  Then $A$ is contained in a generalized arithmetic
progression $P$ of dimension $d \leq d(K)$ and size at most $f(K)\abs{A}$,
where $d(K)$ and $f(K)$ depend only on~$K$.
\end{theorem}

\begin{remark}
Freiman's original (1973) proof gave enormous bounds.  Ruzsa (1994)
gave an elegant proof with $d \leq K - 1$ and $f(K) = e^{CK}$ for some
absolute constant $C > 0$.  The Polynomial Freiman--Ruzsa conjecture
(Gowers 2007, Green--Tao 2006) seeks to sharpen these to
$d = O(\log K)$ and $f(K) = K^{O(1)}$; it was proved by
Gowers--Green--Manners--Tao in 2023.
\end{remark}

\begin{corollary}[Freiman, Double-Sumset Version]
If $\abs{A + A} \leq K\abs{A}$, then $A$ is ``almost'' an arithmetic
progression: it lies in a progression of bounded dimension.
\end{corollary}

\subsection{Cauchy--Davenport Theorem}

For sets in $\Z/p\Z$ ($p$ prime), the analogue of the trivial bound is sharp:

\begin{theorem}[Cauchy--Davenport, \cite{Cauchy1813}]
Let $p$ be prime and $A, B \subseteq \Z/p\Z$ be non-empty.  Then
\[
  \abs{A + B} \;\geq\; \min(p,\; \abs{A} + \abs{B} - 1).
\]
\end{theorem}

This is tight: equality holds when $A$ and $B$ are APs with the same
common difference.

% ============================================================
\section{Erd\H{o}s--Tur\'{a}n Conjecture}\label{sec:erdos-turan}
% ============================================================

\subsection{Representation Functions}

\begin{definition}[Representation Function]
For $A \subseteq \N_0$ and $h \geq 2$, define
\[
  r_{h,A}(n) \;:=\; \abs{\{(a_1, \ldots, a_h) \in A^h :
                          a_1 + \cdots + a_h = n\}}.
\]
If $A$ is a basis of order $h$ for $\N$, then $r_{h,A}(n) \geq 1$
for all sufficiently large~$n$.
\end{definition}

\subsection{The Conjecture}

\begin{conjecture}[Erd\H{o}s--Tur\'{a}n, 1941 \cite{ErdosTuran1941}]
\label{conj:erdos-turan}
If $A \subseteq \N_0$ is an additive basis of order $2$ (i.e., every
sufficiently large integer is the sum of two elements of $A$), then
\[
  r_{2,A}(n) \;=\; \abs{\{(a, b) \in A \times A : a + b = n\}}
\]
cannot be bounded: $\limsup_{n \to \infty} r_{2,A}(n) = \infty$.
\end{conjecture}

In other words, if every large $n$ can be represented as $a + b$ with
$a, b \in A$, then necessarily infinitely many integers have very
\emph{many} such representations.

\begin{remark}
This conjecture is \textbf{open}.  It has been verified for many
specific families of sets, but no general proof is known.

Partial results:
\begin{itemize}
\item Erd\H{o}s (1956): There exists a basis $A$ of order 2 for which
      $r_{2,A}(n) = O(\log n)$.  So the conjecture, if true, would be
      essentially tight.
\item Borwein--Choi--Chu (2006): Verified for specific polynomial
      sequences.
\item The analogous statement for order $h \geq 3$ is also open.
\end{itemize}
\end{remark}

\subsection{Connection to the Goldbach Conjecture}

If Goldbach's binary conjecture (Conjecture~\ref{conj:goldbach-binary})
holds, then $\mathbb{P}$ is a basis of order 2 for even integers.
The Erd\H{o}s--Tur\'{a}n conjecture would then predict
$\limsup_{n \to \infty} G(n) = \infty$, which is consistent with
(but much weaker than) the Hardy--Littlewood conjecture
(Conjecture~\ref{conj:hardy-littlewood}).

% ============================================================
\section{Computational Aspects}\label{sec:computational}
% ============================================================

\subsection{Implementation in \texttt{goldbach\_extended.py}}

The Python module \texttt{goldbach\_extended.py} implements:
\begin{itemize}
\item \texttt{GoldbachExtended}: computes $G(n)$, the Goldbach comet,
      the Hardy--Littlewood heuristic, Chen's theorem verification,
      ternary Goldbach checks, and the Schnirelmann constant analysis.
\item \texttt{VinogradovMethod}: computes the exponential sum $S(\alpha)$,
      the singular series $\mathfrak{S}(N)$, major arc descriptions, and
      a numerical estimate of $r_3(N)$ via numerical integration.
\end{itemize}

\subsection{Goldbach Comet}

The \emph{Goldbach comet} is the plot of $G(n)$ for even $n$.
Table~\ref{tab:comet} shows computed values:

\begin{table}[h]
\centering
\begin{tabular}{@{}cc|cc|cc@{}}
\toprule
$n$ & $G(n)$ & $n$ & $G(n)$ & $n$ & $G(n)$ \\
\midrule
4   & 1  & 20  & 4  & 60  & 9  \\
6   & 1  & 30  & 5  & 100 & 6  \\
8   & 1  & 40  & 3  & 120 & 12 \\
10  & 2  & 50  & 4  & 200 & 9  \\
12  & 1  & 52  & 3  & 360 & 22 \\
\bottomrule
\end{tabular}
\caption{Goldbach partition function $G(n)$ for selected even~$n$.
The spike at $n = 360$ reflects its many small prime factors.}
\label{tab:comet}
\end{table}

\subsection{Hardy--Littlewood Heuristic}

The \texttt{hardy\_littlewood\_goldbach(n)} method computes:
\[
  \widehat{G}(n) \;=\; 2\,\Pi_2
  \prod_{\substack{p > 2 \\ p \mid n}} \frac{p-1}{p-2}
  \cdot \frac{n}{(\log n)^2}.
\]

\begin{example}
For $n = 100$: $100 = 2^2 \cdot 5^2$, so $p = 5$ is the only relevant
odd prime factor.
\[
  \widehat{G}(100) \;=\; 2 \times 0.6602 \times \frac{4}{3} \times
  \frac{100}{(\log 100)^2} \;\approx\; 6.5.
\]
The true value is $G(100) = 6$, a very good fit.
\end{example}

\subsection{Ternary Goldbach Verification}

The method \texttt{verify\_ternary\_goldbach\_range(n\_max)} confirms
Helfgott's theorem computationally for $n \leq 10{,}000$:
\begin{itemize}
\item All 4{,}997 odd numbers from 7 to 9{,}999 can be expressed as
      a sum of three primes.
\item No failures found (consistent with Theorem~\ref{thm:helfgott}).
\end{itemize}

\subsection{Singular Series Computation}

For odd $N = 101$, the singular series (computed with \texttt{prime\_limit=200}) gives:
\[
  \mathfrak{S}(101) \;\approx\; 1.3284.
\]
The predicted number of representations (weighted) is then:
\[
  r_3(101) \;\approx\; 1.3284 \times \frac{101^2}{2 (\log 101)^3}
  \;\approx\; 74.
\]

\subsection{Waring's Problem --- Verification}

The module \texttt{additive\_number\_theory.py} implements
\texttt{WaringProblem.waring\_g(k)} which verifies $g(k)$ for small~$k$
by exhaustive search:
\[
  g(2) = 4,\quad g(3) = 9,\quad g(4) = 19.
\]
These are confirmed by direct computation for integers up to
$n = 1{,}000$.

% ============================================================
\section{Summary and Open Problems}\label{sec:summary}
% ============================================================

We summarize the status of the main results:

\begin{table}[h]
\centering
\begin{tabular}{@{}p{6.5cm}ll@{}}
\toprule
Statement & Status & Reference \\
\midrule
Binary Goldbach ($n \geq 4$ even $= p + q$) &
  \textbf{Conjecture (open)} & \cite{Goldbach1742} \\
Ternary Goldbach ($n \geq 7$ odd $= p+q+r$) &
  \textbf{Theorem} & \cite{Helfgott2013} \\
Chen's Theorem (large even $n = p + P_2$) &
  \textbf{Theorem} & \cite{Chen1973} \\
Vinogradov's 3-prime theorem (large odd) &
  \textbf{Theorem} & \cite{Vinogradov1937} \\
Hardy--Littlewood Conjecture ($G(n)$ asymptotic) &
  \textbf{Conjecture (open)} & \cite{HardyLittlewood1923} \\
Hilbert--Waring Theorem ($g(k) < \infty$) &
  \textbf{Theorem} & \cite{Hilbert1909} \\
Lagrange Four Squares ($g(2) = 4$) &
  \textbf{Theorem} & \cite{Lagrange1770} \\
Schnirelmann: primes form a basis &
  \textbf{Theorem} & \cite{Schnirelmann1930} \\
Freiman's Structure Theorem &
  \textbf{Theorem} & \cite{Freiman1973} \\
Erd\H{o}s--Tur\'{a}n Conjecture &
  \textbf{Conjecture (open)} & \cite{ErdosTuran1941} \\
\bottomrule
\end{tabular}
\caption{Overview: proved theorems vs.\ open conjectures.}
\label{tab:status}
\end{table}

\subsection*{Open Problems}
\begin{enumerate}[label=(\arabic*)]
\item Prove (or disprove) the binary Goldbach conjecture.
\item Determine the exact value of $G(k)$ for $k \geq 3$.
\item Prove the Erd\H{o}s--Tur\'{a}n conjecture.
\item Prove the Hardy--Littlewood conjecture on $G(n)$.
\item Extend Green--Tao (primes contain arbitrarily long APs) to
      quantitative results on the length of APs in dense subsets of primes.
\end{enumerate}

% ============================================================
\section*{References}\label{sec:references}
% ============================================================

\begin{thebibliography}{99}

\bibitem{Goldbach1742}
C.\ Goldbach,
\textit{Letter to L.\ Euler}, 7 June 1742.

\bibitem{Lagrange1770}
J.-L.\ Lagrange,
\textit{D\'{e}monstration d'un th\'{e}or\`{e}me d'arithm\'{e}tique},
Nouveaux M\'{e}m.\ Acad.\ Roy.\ Sci.\ Belles-Lettres de Berlin (1770), 123--133.

\bibitem{Cauchy1813}
A.-L.\ Cauchy,
\textit{Recherches sur les nombres},
J.\ \'{E}cole Polytech.\ \textbf{9} (1813), 99--116.

\bibitem{Schnirelmann1930}
L.\ G.\ Schnirelmann,
\textit{\"{U}ber additive Eigenschaften von Zahlen}
(Russian), Math.\ Ann.\ \textbf{107} (1932), 649--690.
Originally in: Izv.\ Donskogo Politekhn.\ Inst.\ \textbf{14} (1930).

\bibitem{Vinogradov1937}
I.\ M.\ Vinogradov,
\textit{Some theorems concerning the theory of prime numbers},
Matem.\ Sbornik \textbf{2}(44) (1937), 179--196.

\bibitem{HardyLittlewood1923}
G.\ H.\ Hardy and J.\ E.\ Littlewood,
\textit{Some problems of `Partitio Numerorum': III.\ On the expression of
a number as a sum of primes},
Acta Math.\ \textbf{44} (1923), 1--70.

\bibitem{Hilbert1909}
D.\ Hilbert,
\textit{Beweis f\"{u}r die Darstellbarkeit der ganzen Zahlen durch eine
feste Anzahl $n^{\text{ter}}$ Potenzen (Waringsche Problem)},
Math.\ Ann.\ \textbf{67} (1909), 281--300.

\bibitem{Chen1973}
J.\ R.\ Chen,
\textit{On the representation of a large even integer as the sum of a
prime and the product of at most two primes},
Sci.\ Sinica \textbf{16} (1973), 157--176.

\bibitem{ErdosTuran1941}
P.\ Erd\H{o}s and P.\ Tur\'{a}n,
\textit{On a problem of Sidon in additive number theory, and on some
related problems},
J.\ London Math.\ Soc.\ \textbf{16} (1941), 212--215.

\bibitem{Freiman1973}
G.\ A.\ Freiman,
\textit{Foundations of a Structural Theory of Set Addition},
Transl.\ Math.\ Monographs \textbf{37}, AMS, Providence, RI, 1973.

\bibitem{OliveiraSilva2012}
T.\ Oliveira e Silva, S.\ Herzog, and S.\ Pardi,
\textit{Empirical verification of the even Goldbach conjecture and
computation of prime gaps up to $4 \times 10^{18}$},
Math.\ Comp.\ \textbf{83} (2014), 2033--2060.

\bibitem{Helfgott2013}
H.\ A.\ Helfgott,
\textit{The ternary Goldbach conjecture is true},
arXiv:1312.7748 (2013).

\end{thebibliography}

\end{document}
